\documentclass[12pt,a4paper]{article}
\usepackage{graphicx}
\usepackage{amsmath}
\usepackage{hyperref}
\usepackage{geometry}
\geometry{margin=2.5cm}
\emergencystretch=3em

\title{\textbf{CDFD Runtime Paper 02: Layered Ontology for Adaptive Systems}}
\author{Steve Bico Mujjabi, MD\\
\small Independent Researcher\\
\small ORCID: \href{https://orcid.org/0009-0001-0556-5516}{0009-0001-0556-5516}}
\date{May 2026}

\begin{document}
\maketitle

\begin{abstract}
\noindent \textbf{Executive synthesis.} The CDFD Runtime separates numerical
mechanics from meaning. The engine evolves arrays; the ontology layer explains
what those arrays represent in a domain. This paper upgrades the earlier
``5-layer'' account into a more explicit layered stack: runtime mechanics,
meta-ontology, domain ontologies, application ontologies, and context-specific
memory. It also explains how the discovery layer can propose new schemas, while
keeping every new schema in candidate status until it survives validation.
\end{abstract}

\paragraph{Greek-symbol convention.}
Unless a manuscript states a tighter equation-local meaning, Greek symbols are local CDFD variables or coefficients, not universal constants. Here $\Phi$ denotes driving flux, $\Psi_s$ denotes the adaptive operating ratio, $\Omega$ denotes a domain or state space, and $\Lambda$ denotes the Life Number where used. The coefficients $\alpha$, $\beta$, and $\gamma$ denote accumulation or reinforcement, relaxation or decay, and diffusion, coupling, or redistribution. The symbols $\delta$ and $\epsilon$ denote perturbation or tolerance scales, while $\eta$, $\lambda$, $\mu$, $\nu$, $\rho$, $\sigma$, $\tau$, and $\phi$ denote local efficiency, rate, baseline, frequency, density or correlation, noise or variance, time-scale, and phase or field terms. Subscripts restrict any coefficient to the named pathway, tissue, domain, or experiment.


\begin{figure}[ht]
\centering
\fbox{\begin{minipage}{0.92\textwidth}
\centering
\textbf{Runtime evidence path}\\[0.4em]
Prior CDFD mathematics $\longrightarrow$ CDFL/runtime state $\longrightarrow$ bounded output $\longrightarrow$ falsification target.
\end{minipage}}
\caption{Conceptual schematic for the runtime papers. The engine translates the mathematical frame from the prior CDFD papers into executable CDFL/runtime diagnostics; candidate outputs remain tied to explicit tests before any stronger claim is made.}
\end{figure}

\section{Prior CDFD Basis}
This manuscript does not rederive the mathematical core of CDFD. It uses the
Mujjabi Adaptive Operating Ratio and the constrained-vacuum notation introduced
in Part I \cite{MujjabiI2026}, the torus-knot hierarchy and boundary-mode work
from the Part I sequence \cite{MujjabiVI2026,MujjabiIX2026}, the public balance
law developed in the Part I capstone \cite{MujjabiX2026}, the Part I transport
tests \cite{MujjabiXII2026}, and the Life Number and tri-regime biology bridge
from Part II \cite{MujjabiPartII2026}. The runtime papers therefore act as an
execution layer: they keep the mathematics connected to commands, outputs,
falsification tests, and domain-specific limits.


\section{Why Ontology Is Needed}
The core engine does not know what a nephron, a market, a mineral pore, or a
photonic channel is. It only evolves state fields such as $\Phi$, $C$, $S$,
$M_s$, and $\Psi_s$. That ignorance is deliberate. A reusable runtime must not
hard-code every domain into the kernel.

The ontology layer supplies meaning without contaminating the engine. It maps
domain concepts to runtime variables, records relationships, and exposes those
relationships to CDFL, the CLI, and the future web interface.

\section{The Runtime Ontology Stack}
\subsection{Layer 0: Runtime Mechanics}
The runtime layer is responsible for graph mechanics, state propagation,
uncertainty flags, causal links, memory binding, regime tagging, and finite
diagnostic summaries. It does not decide what a domain-specific state means.

\subsection{Layer 1: Core Meta-Ontology}
The meta-ontology defines reusable categories:
\begin{itemize}
    \item \textbf{Entity}: a node or bounded system.
    \item \textbf{Process}: a change pathway.
    \item \textbf{Flow}: a domain expression of $\Phi$.
    \item \textbf{Constraint}: a domain expression of $C$.
    \item \textbf{Memory}: a domain expression of $M_s$.
    \item \textbf{Regime}: a qualitative label derived from runtime state.
\end{itemize}

\subsection{Layer 2: Domain Ontologies}
Domain adapters in \texttt{cdfd\_runtime/domains} translate domain records into
engine variables. The CLI currently reports 196 registered domain keys. The
number matters less than the contract: each adapter must define how a domain
maps to $\Phi$ and $C$, and must explain its interpretation limits.

\subsection{Layer 3: Application Ontologies}
Applications such as CDFD Studio, future scientific dashboards, or domain tools
call the shared runtime backend rather than invent their own physics.
This preserves one source of truth for finite checks, result envelopes, and
CDFL execution.

\subsection{Layer 4: Context and Memory}
Context ontologies store user, project, dataset, or experiment-specific
relationships. These guide interpretation, but they do not
silently alter the physics kernel.

\section{Schema Invention and Discovery}
The local report \texttt{experiments/reports/NEW\_DISCOVERIES\_REPORT.md}
describes a schema-invention demonstration in which a high-$\Psi_s$ anomaly was
assigned a new candidate label. The proper interpretation is not that the
runtime has proven a new law of nature. The useful result is narrower: the
discovery layer can detect an out-of-distribution regime, name it, record its
features, and create a target for future falsification.

\section{CDFL and Ontology Binding}
CDFL binds ontology to execution. A \texttt{SET domain: physics} directive
selects a semantic context; a \texttt{SYSTEM} block maps flow and constraint;
a \texttt{RULE} block attaches an action to a threshold. The CLI now exposes
this pipeline through:
\begin{verbatim}
python cdfd.py validate examples/heat_flow.cdfl
python cdfd.py run examples/heat_flow.cdfl
\end{verbatim}

\section{Falsification Conditions}
An ontology mapping fails if:
\begin{itemize}
    \item it cannot define measurable proxies for $\Phi$ and $C$;
    \item its predicted regime thresholds do not reproduce under rerun;
    \item it adds no explanatory or predictive value beyond existing variables;
    \item it hides uncertainty behind authoritative labels.
\end{itemize}

\section{Conclusion}
The ontology is not decoration. It is the meaning infrastructure that lets one
runtime serve many domains without pretending those domains are materially
identical. Its strongest use is disciplined translation: map, run, report,
falsify, and only then promote a candidate schema.
\begin{thebibliography}{9}
\bibitem{MujjabiI2026}
Steve Bico Mujjabi, MD. \emph{Vortex Stability in a Constrained Transport Vacuum and the Origin of the Fine-Structure Constant}. CDFD Part I Paper I, preprint, 2026.

\bibitem{MujjabiVI2026}
Steve Bico Mujjabi, MD. \emph{The Universal Torus Knot Hierarchy: $Z_n$ Symmetry, the Cinquefoil Family, and the General Structure of CDFT Particle Families}. CDFD Part I Paper VI, preprint, 2026.

\bibitem{MujjabiIX2026}
Steve Bico Mujjabi, MD. \emph{Even-$n$ Torus Knots in CDFT: The Exclusion of $T(2,2)$, the Extension to $T(2,2m)$, and the Anti-Phase Mass Scaling Law}. CDFD Part I Paper IX, preprint, 2026.

\bibitem{MujjabiX2026}
Steve Bico Mujjabi, MD. \emph{The Vacuum Equation of State: Deriving Torus Knot Self-Consistency from the Public CDFD Balance Law $\Psi = \Phi/C$}. CDFD Part I Paper X, preprint, 2026.

\bibitem{MujjabiXII2026}
Steve Bico Mujjabi, MD. \emph{Friction, Superconductivity, Plasma States, and Transport Mysteries: Observable Tests of Adaptive Constraint}. CDFD Part I Paper XII, preprint, 2026.

\bibitem{MujjabiPartII2026}
Steve Bico Mujjabi, MD. \emph{CDFD Part II: Origins of Life and Tri-Regime Bioenergetics}. Public release archive, 2026, \url{https://doi.org/10.5281/zenodo.20264779}.
\end{thebibliography}
\end{document}
